%% kapa: exam=1 round=28 year=2017 period=1 form=A subject=civil_law subject_ko=민법 seq=40 number=40
%% source: Q-Net 감정평가사 기출문제 (cst00309) artlSeq=5199603 file=extracted/1008/5199603/2017년 제28회 감정평가사 1차시험 문제지/1교시 A형.pdf
%% conversion: kordoc 4.24.51 -> tools/gen1.py (auto); review=proofread-sonnet; flags=none
\begin{question}{40}
물권적 청구권에 관한 설명으로 옳지 않은 것은? (다툼이 있으면 판례에 따름)
\begin{choices}
  \item 소유권에 기한 물권적 청구권은 소멸시효 대상이 아니다.
  \item 간접점유자는 제3자의 점유침해에 대하여 물권적 청구권을 행사할 수 있다.
  \item 토지의 저당권자는 무단점유자에 대해 저당권에 기한 저당물반환청구권을 행사할 수 있다.
  \item 점유물반환청구권은 점유를 침탈당한 날로부터 1년 내에 행사하여야 한다.
  \item 점유보조자에게는 점유보호청구권이 인정되지 않는다.
\end{choices}
\end{question}
\answer{3}
